% SPDX-License-Identifier: LPPL-1.3c
% Copyright (C) 2026 Christophe Jorssen
% Author and Current Maintainer: Christophe Jorssen <christophe.jorssen@gmail.com>
% LPPL maintenance status: maintained
% This file is part of luafemm. See LICENSE and LICENSES.md for its terms.
\section{Saving physical problems and answers}
\label{sec:native-storage}

Keep the original TikZ source for editing the drawing. A |.femx| document stores
the evaluated physical problem; an |.ansx| document additionally owns the answer.
These are \textsf{luafemm} formats. FEMM does not read them: use the separate
|.fem|/|.ans| or |.fee|/|.res| conversion interfaces when exchanging with FEMM.
Unlike a cache, an answer is selected explicitly and owns its physical data.
It does not inherit the current material catalogue or silently recompute a field.

\begin{command}{\femmsave\oarg{options}}
Save the current model at the point of invocation. Declare all physical paths
and sources first. The command groups its |/femm/save| options. An answer save
solves an unsolved builder; a problem save never solves. Validation precedes
atomic replacement, so an unsupported model leaves an existing target intact.
\end{command}
\begin{key}{/femm/save/file=\meta{filename} (initially empty)}
Required output filename. Use |.femx| for a problem, |.ansx| for an answer.
Relative outputs follow the engine's output directory. Parent directories must
already exist. File contents are passive data, not executable Lua or TeX.
\end{key}
\begin{key}{/femm/save/kind=problem\textbar answer (initially problem)}
Select physical declarations alone or a physical problem with its solution.
\end{key}
\begin{key}{/femm/save/physics=magnetism\textbar electrostatics (initially magnetism)}
Select the model registry explicitly; an electric picture does not change this
default. This makes mixed-physics documents independent of drawing order.
\end{key}
\begin{key}{/femm/save/model=\meta{name} (initially empty)}
An empty value uses the current picture. Otherwise select a named magnetic
model or electric snapshot. Electric snapshots have already been frozen.
\end{key}
\begin{key}{/tikz/femm/save=\meta{options}}
Invoke the same operation as |\femmsave| through a TikZ key. On a path, put it
after all physical declarations have been captured; a region's postaction has
not yet run while its own options are being parsed. Prefer a separate
|\tikzset{femm/save={...}}| at the end of the physical declarations.
\end{key}
\begin{key}{/femm/problem/native file=\meta{filename} (initially empty)}
Read a magnetic |.femx| or |.ansx| at picture initialization. The file owns the
physical and numerical settings. Do not combine this selector with FEMM import,
an operating point, a sliding template, a time state or explicit problem settings.
Ordinary TikZ transforms and field styles remain available.
\end{key}
\begin{key}{/femm/electrostatic/native file=\meta{filename} (initially empty)}
The electric counterpart. Do not combine it with |file|, |snapshot| or explicit
physical geometry/medium selectors. Stored model units govern all subsequent
coordinates. Restoring an answer reconstructs element gradients and reactions
from its exact ordered mesh and coefficients; it performs no linear or nonlinear
solve. Changing drawing styles never modifies the stored physical problem.
\end{key}

\subsection{One answer, two presentations}
The first picture declares and saves a dipole. The second reads the same answer
and displays potential contours. No material or charge declaration is repeated.
\begin{codeexample}[width=5cm]
\begin{tikzpicture}[scale=.7,femm/electrostatic]
 \node[femm/point charge={charge=1e-9}] at (-1,0) {$+$};
 \node[femm/point charge={charge=-1e-9}] at (1,0) {$-$};
 \femmsave[physics=electrostatics,kind=answer,
   file=manual-dipole.ansx]
 \pic[femm/electric={lines per charge=8}]
   {femm electric field};
\end{tikzpicture}
\end{codeexample}
\begin{codeexample}[width=5cm]
\begin{tikzpicture}[scale=.7,
 femm/electrostatic={native file=manual-dipole.ansx}]
 \pic[femm/electric={levels={-900,-450,0,450,900}}]
   {femm equipotentials};
 \node at (-1,0) {$+$}; \node at (1,0) {$-$};
\end{tikzpicture}
\end{codeexample}

\subsection{Representations and limits}
Analytical charge and symmetry providers store registered law parameters;
they do not invent a mesh. Electric planar and axisymmetric P1 answers retain
the scalar nodal potential. Source-correction answers retain their singular
kernel parameters and the regular correction coefficients separately. Thus
their finite interaction energy and total-field contours remain available.
Ordinary static and harmonic magnetic planar/axisymmetric answers retain
the ordered P1 vector-potential coefficients, with separate real and imaginary
arrays for peak |exp(+i omega t)| phasors. Embedded material values are authoritative.
The prescribed-current 3D adapter stores its registered source laws and quadrature
settings without a fabricated mesh; see Section~\ref{sec:biot}.

The adapters exclude imported electric native topology, ideal magnetic circuits,
axial diffusion, operating responses, sliding bands, time states, network and
flux-controlled magnetic models, and Lua callbacks. Unsupported adapters fail;
they are never converted to a simpler physical model. An answer is not a
transient restart: use the separately supported |.lft| state interface.
The |.lfc| cache contract also remains independent.

\subsection{Versioned passive records}
The text header is |LUAFEMM-FEMX| or |LUAFEMM-ANSX|, followed by payload byte count
and an Adler-32 checksum. The bounded codec stores finite numbers, booleans,
length-prefixed strings and sorted maps. It never calls |load| on input.
Schema version 1 declares physics, physical dimension, evaluation dimension (two for section providers, three for prescribed currents),
coordinate scale, SI field units, potential reference, required adapter and
solution basis. The physical dimension is three even for a planar section;
planar invariance and axisymmetry remain distinct geometry tags.
Package provenance is descriptive; schema and adapter revisions decide
compatibility. Unknown required capabilities and future schema versions fail
before reconstruction. No implicit migration, external dependency loading or
original TikZ macro reconstruction is performed. The checksum detects accidental
damage, not malicious modification or scientific invalidity.
